% Chapter 1: Introduction

\section{Background}

Botswana's agricultural sector is more important than people realize, especially livestock farming. The cattle population exceeds 2 million head, which makes it one of the largest per capita in Africa \cite{botswana_agriculture_2023}. That's a lot of cattle, and managing them the traditional way is becoming increasingly difficult.

The problem is that farmers still rely on manual observation and periodic physical inspections. This is labor-intensive and, more importantly, it means health issues are often detected too late. By the time symptoms are visible, treatment is less effective and mortality rates increase. I've seen this happen firsthand - farmers lose cattle that could have been saved if the problem had been caught earlier.

IoT technologies could change this. Sensor networks can provide continuous, real-time monitoring of cattle health, enabling early detection and timely intervention. But implementing IoT in agricultural contexts isn't straightforward. The data governance challenges are significant - how do you ensure data quality, protect farmer privacy, and maintain system reliability when you're dealing with resource-constrained environments and intermittent connectivity?

\section{Problem Statement}

The current livestock management practices in Botswana have several critical problems that need to be addressed:

1. **Delayed Health Detection**: Without real-time monitoring, health issues are often detected only after symptoms become severe. This reduces treatment effectiveness and increases mortality rates. I've spoken to farmers who lost entire herds to diseases that could have been contained if detected earlier.

2. **Limited Data Governance**: Existing agricultural systems lack robust data governance frameworks. This leads to inconsistent data quality, security vulnerabilities, and privacy concerns. During my research, I found that many systems don't even have basic validation rules - they accept whatever data comes in, regardless of whether it makes sense.

3. **Poor Data Integrity**: Many agricultural IoT systems don't implement proper database normalization. This results in data redundancy, update anomalies, and maintenance challenges. When I looked at existing systems, I found that they often store the same data in multiple places, which makes updates a nightmare and leads to inconsistencies.

4. **Privacy Concerns**: Farmer data, including personal information and location data, is often stored without adequate privacy protections. This raises concerns about data misuse and unauthorized access. Farmers I spoke to were particularly worried about competitors accessing their location data and cattle counts.

5. **Scalability Issues**: Traditional systems cannot scale to handle the volume of telemetry data generated by large-scale IoT deployments. This limits their applicability for commercial operations. A system that works for 50 cattle breaks when you try to scale to 500.

These challenges highlight the need for a comprehensive data governance framework that addresses database design, data integrity, privacy protection, and system scalability specifically tailored for agricultural IoT contexts in developing regions like Botswana.

\section{Motivation}

I chose this project for both practical and academic reasons:

**Practical Motivation:**
- Economic losses due to livestock diseases in Botswana are estimated at millions of Pula annually \cite{botswana_livestock_2022}. This is real money that farmers can't afford to lose.
- Improved livestock health monitoring could significantly reduce these losses. I've seen farmers struggle with disease outbreaks, and earlier detection would make a real difference.
- Technology adoption in Botswana's agricultural sector is increasing. Farmers are becoming more comfortable with digital tools, which creates an opportunity for IoT solutions.
- Privacy concerns are becoming increasingly important as digital systems become more prevalent. Farmers are asking questions about who can access their data, and they deserve answers.

**Academic Motivation:**
- There's limited research on data governance frameworks specifically designed for agricultural IoT in developing regions. Most research assumes developed-country infrastructure.
- The intersection of database normalization, privacy-by-design, and agricultural IoT presents an interesting research area. It's not just about applying existing concepts - it's about adapting them for a specific context.
- This project gives me the opportunity to apply theoretical database concepts (3NF normalization) in a practical context. Database theory is abstract, but seeing it work in an actual system is different.
- The project aligns with INFS 401 module requirements for data governance and database design, which is important for my degree.

\section{Aim}

My aim is to design and implement a comprehensive data governance framework for agricultural IoT systems in Botswana. The focus is on database design, data integrity, privacy protection, and system scalability. I want to build something that actually works in the field, not just in a lab.

\section{Objectives}

The specific objectives I set for this project are:

1. Design and implement a Third Normal Form (3NF) database schema for agricultural telemetry data, ensuring data integrity and eliminating redundancy. This is foundational - if the database is poorly designed, everything else fails.

2. Develop a data validation engine that enforces data quality standards and prevents invalid data entry. I want to catch bad data before it ever enters the system.

3. Implement a Role-Based Access Control (RBAC) system to protect sensitive farmer data and ensure privacy-by-design. Farmers need to know that their data is protected.

4. Create a data quality scoring system to monitor and assess the quality of telemetry data over time. This helps identify when sensors are drifting or when data quality is degrading.

5. Validate the system through comprehensive testing to ensure compliance with data governance standards. Testing is what separates theory from practice.

6. Document the system according to academic standards, including normalization proofs, security design, and performance analysis. The documentation needs to be thorough enough that someone else could understand and maintain the system.

\section{Scope}

\subsection{In-Scope}
- 3NF database schema design for agricultural telemetry data
- Data validation engine with quality scoring
- RBAC implementation with permission matrix
- Privacy-by-design features including anonymized farmer IDs
- Database performance optimization with indexing
- Data integrity constraints and foreign key relationships
- Audit logging for data access tracking

These are the core data governance features that I can realistically implement within the INFS 401 scope. I focused on database design, validation, access control, and privacy because these are the foundational elements of any data governance framework.

\subsection{Out-of-Scope}
- Multi-agent orchestration algorithms (covered by COMP 401 project)
- Graph-based routing systems (covered by COMP 401 project)
- Hardware sensor integration
- Mobile application development
- Machine learning models for health prediction
- Blockchain integration for provenance tracking
- Real-time data streaming and processing

I deliberately excluded these features to avoid scope creep and to maintain clear separation between INFS 401 and COMP 401. The multi-agent orchestration and graph routing are complex algorithmic systems that belong in COMP 401, not INFS 401. Hardware integration would require physical sensors and field testing, which isn't feasible within the academic timeline. Machine learning and blockchain would add significant complexity without directly addressing the core data governance problem.

\section{Technologies}

I chose these technologies for specific reasons:

- **Python 3.10+**: Primary programming language for backend implementation. Python is widely used, has excellent database libraries, and is easy to debug. I chose it over alternatives like Java or Go because of the learning resources available and the speed of development.

- **SQLite**: Database management system for edge deployment. SQLite is perfect for this use case because it's serverless, requires no configuration, and works well on resource-constrained devices. I considered PostgreSQL but decided against it because it would require a separate database server, which adds complexity for edge deployment.

- **Pydantic**: Data validation library for schema validation. Pydantic provides type-safe data validation with minimal boilerplate. I initially considered writing custom validation logic, but Pydantic handles the common cases and allows me to focus on the agricultural-specific validation rules.

- **Flask**: Web framework for API layer (minimal implementation). I'm using Flask minimally - just enough to provide an API interface for the COMP project to consume. I didn't need a full web application, so Flask's simplicity was appealing.

- **LaTeX**: Document preparation system for dissertation. LaTeX is the standard for academic writing, and it handles citations and formatting automatically. The learning curve was steep, but it's worth it for the professional output.

\section{Risks}

\subsection{Technical Risks}
- **Database Performance**: SQLite may not scale to handle very large datasets. However, for the expected data volume in this project, this should be acceptable. If performance becomes an issue, I can add more indexes or migrate to PostgreSQL in a future iteration.
- **Data Validation Complexity**: Implementing comprehensive validation rules may be more complex than initially anticipated. Agricultural data has edge cases I haven't encountered yet, so the validation engine may need to be extended as I discover new scenarios.
- **RBAC Implementation**: Designing an appropriate permission matrix requires careful consideration of user roles and access patterns. I need to make sure the roles actually reflect how farmers, administrators, and brokers interact with the system.

\subsection{Project Risks}
- **Time Constraints**: The academic semester timeline may limit the depth of implementation for some features. I've prioritized the core data governance features, but some nice-to-have items may need to be deferred to COMP 402.
- **Learning Curve**: Some technologies (Pydantic, advanced SQLite features) may require additional learning time. I've budgeted time for this, but unexpected complexity could delay other parts of the project.
- **Testing Limitations**: Limited access to real-world agricultural environments may constrain comprehensive testing. I'll need to simulate realistic test scenarios, which may not capture all the edge cases that would occur in actual deployment.

\subsection{Academic Risks}
- **Scope Creep**: The temptation to add features beyond the INFS 401 scope (e.g., multi-agent systems) must be carefully managed. I need to stay focused on data governance and not drift into COMP 401 territory.
- **Documentation Quality**: Balancing technical depth with academic writing standards requires careful attention. The dissertation needs to be academically rigorous while still being technically accurate.
- **Assessment Alignment**: Ensuring the project clearly demonstrates INFS 401 learning outcomes without overlapping with COMP 401. The integration architecture document helps with this, but I need to be vigilant about maintaining clear module boundaries.

\section{Dissertation Structure}

This dissertation is organized as follows:

- **Chapter 1**: Introduction - Provides background, problem statement, motivation, aim, objectives, scope, and risks
- **Chapter 2**: Literature Review - Examines existing research on database normalization, data governance, privacy-by-design, and agricultural IoT
- **Chapter 3**: Methodology - Describes the research design, development approach, and evaluation methods
- **Chapter 4**: Analysis - Presents functional and non-functional requirements, use cases, and UML diagrams
- **Chapter 5**: Design - Details the database schema, validation engine design, RBAC design, and security architecture
- **Chapter 6**: Implementation - Describes the implementation of the data governance framework
- **Chapter 7**: Testing - Presents testing methodology, test cases, and results
- **Chapter 8**: Results - Discusses the performance and effectiveness of the implemented system
- **Chapter 9**: Discussion - Interprets results, compares with literature, and discusses limitations
- **Chapter 10**: Conclusion - Summarizes contributions, identifies limitations, and suggests future work

The structure follows the standard academic format for a technical dissertation, with each chapter building on the previous one. I've tried to make the narrative flow naturally from problem identification through solution implementation to evaluation and reflection.
